\chapter{Many machines: exploration, regret, and the first bonus term}
\label{ch:manymachines}

\begin{objectives}
\begin{itemize}
\item Define regret, and use it to say precisely what "wasted thinking" means.
\item Show by explicit counterexample that always picking the best-so-far move is a
  catastrophic policy, and that adding randomness is a poor repair.
\item Build the UCB1 rule from Chapter~\ref{ch:onemachine}'s confidence bound, through
  three failed attempts, and understand why the final bonus contains $\ln N$ over $n_i$
  under a square root.
\item Run UCB1 by hand for twelve pulls over three machines and watch it self-correct.
\item State what UCB1 guarantees --- and where that guarantee stops applying to chess.
\end{itemize}
\end{objectives}

\section{The setup, and the honest definition of waste}

You face $K$ machines. Machine $i$ has an unknown true value $\mu_i$; pulling it returns a
noisy sample. You have $T$ pulls in total. In chess: $K$ is the number of legal moves at the
node ($\approx 31$, Chapter~\ref{ch:problem}), a pull is one visit, and $T$ is the node
budget --- 800, or 1{,}600, or whatever the tool allows.

How do we score a strategy? Not by "did it find the best machine", which is too crude. The
standard measure is \textbf{regret}: how much worse you did than a genie who knew $\mu$ from
the start and pulled the best machine every time.

\begin{equation}
  \label{eq:regret}
  \text{Regret}(T) \;=\; \underbrace{T\,\mu^{*}}_{\text{the genie's total}}
  \;-\; \underbrace{\sum_{t=1}^{T} \mu_{i_t}}_{\text{your total}},
  \qquad \mu^{*} = \max_i \mu_i .
\end{equation}

Equivalently, if machine $i$ is worse than the best by $\Delta_i = \mu^{*} - \mu_i$ and you
pulled it $n_i$ times, then
\begin{equation}
  \label{eq:regret2}
  \text{Regret}(T) \;=\; \sum_{i} n_i \Delta_i .
\end{equation}

Equation~\eqref{eq:regret2} is the one to remember, because it says something a chess player
recognises: \textbf{a bad move is not expensive to look at if it is only slightly bad, and
looking briefly at a terrible move costs almost nothing --- what ruins you is spending many
visits on something clearly inferior.} A search that spends 700 of its 800 nodes on a move
that loses is not "thorough", it is a search with enormous regret.

\begin{notationbox}
$K$ machines/moves; $\mu_i$ true value of $i$; $\mu^*$ the best; $\Delta_i = \mu^* - \mu_i$
its \emph{gap}; $n_i$ times pulled; $N = \sum_i n_i$ total pulls so far; $Q_i$ the running
average of machine $i$ (Chapter~\ref{ch:onemachine}).
\end{notationbox}

\section{Attempt 0: just pick the best so far (greedy)}

The obvious rule: pull $\argmax_i Q_i$.

\begin{worked}{how greedy destroys itself}
Two machines. Truth (unknown to the algorithm): $\mu_A = 0.9$, $\mu_B = 0.5$. Samples are
noisy, so any single pull can mislead.

\begin{enumerate}
\item Pull each once. $A$ has an unlucky sample: $x_A = 0.1$. $B$ has a fair one:
  $x_B = 0.5$. Now $Q_A = 0.1$, $Q_B = 0.5$.
\item Greedy pulls $\argmax Q = B$ forever. Every $B$ sample keeps $Q_B$ near $0.5 > 0.1$.
  $A$ is never touched again, so $Q_A$ \emph{never changes}.
\end{enumerate}
Over $T = 800$ pulls the regret is roughly $798 \times (0.9-0.5) = 319$: essentially the
maximum possible. And note the mechanism --- \textbf{the algorithm has no way of ever
discovering its error}, because the only thing that could correct it is the very action it
has stopped taking.
\end{worked}

That failure has a precise chess name. A single shallow look at a strong move can return a
bad number (the refutation of the refutation is two plies further on). If the search then
never returns to that move, the engine has locked in a blunder \emph{and cannot detect it}.
Any usable search must guarantee that a mistakenly-dismissed move gets another chance.

\subsection{Attempt 0b: greedy with a coin ($\epsilon$-greedy)}

Standard first repair: with probability $\epsilon$ (say $0.1$) pull a machine at random,
otherwise be greedy. This does fix the fatal flaw --- $A$ will eventually be re-sampled ---
but it is crude in two ways that matter for chess.

\begin{enumerate}
\item \textbf{It never stops wasting.} It keeps spending $10\%$ of the budget uniformly
  at random forever, including on machines already known to be terrible. In chess, with
  $K\approx31$, that means a permanent tax paid mostly on moves that hang a piece.
\item \textbf{It explores blindly.} It cannot distinguish "unpromising and well-measured"
  from "unpromising and barely looked at". Those deserve completely different treatment,
  and $\epsilon$-greedy treats them identically.
\end{enumerate}

We already have the tool that fixes both: the confidence bound of
Chapter~\ref{ch:onemachine}. Exploration should be aimed at \emph{uncertainty}, not sprayed.

\section{Building the bonus, one failure at a time}

We want a rule of the form
\begin{equation}
  \label{eq:scoreform}
  \text{score}_i \;=\; Q_i \;+\; (\text{bonus}_i),
\end{equation}
pull the highest score, where the bonus expresses "how much room there might be above what
I have measured". Chapter~\ref{ch:onemachine} says the room is $\epsilon = R\sqrt{\ln(1/\delta)/2n_i}$.
The only real question is what to put in place of $\delta$. Let us get there by failing
three times.

\begin{attempt}{1}{a constant bonus}
\[
\text{score}_i = Q_i + c .
\]
Useless: adding the same constant to everyone changes nothing about the ranking. The bonus
must depend on how much you have looked at each machine.
\end{attempt}

\begin{attempt}{2}{bonus $= c/n_i$}
Now the bonus shrinks with evidence, which is the right direction. But it shrinks
\emph{too fast}, and something worse: it is unaffected by the passage of time.

Concretely, with $c = 1$: a machine pulled 10 times has bonus $0.1$. Suppose its $Q$ is
$0.30$ while the current leader sits at $Q = 0.55$. Its score is $0.40$, it never wins the
argmax --- and as the leader is pulled another thousand times, \emph{nothing about our
machine changes}. Its bonus stays $0.1$ forever.

Why is that wrong? Because the meaning of "unexplored" is relative. A machine with 10 pulls
is well understood when the total budget spent is 20, and barely touched when the total is
100{,}000. \textbf{The bonus must grow as the rest of the search grows.}
\end{attempt}

\begin{attempt}{3}{bonus $= c/\sqrt{n_i}$}
Better --- this is the honest $1/\sqrt{n}$ shape of Equation~\eqref{eq:se} --- but it has
exactly the same defect: no dependence on total pulls $N$. A machine can be permanently
frozen out by a rival that got lucky early, no matter how long the search runs.

We need a numerator that grows with $N$.
\end{attempt}

\subsection{Choosing $\delta$: how the $\ln N$ appears}

Return to Equation~\eqref{eq:epsilon}: $\epsilon = R\sqrt{\ln(1/\delta)/2n_i}$. We must
choose the failure probability $\delta$. Here is the reasoning that fixes it.

We are not making one statement; we are making one statement \emph{per machine, per time
step}, and we would like all of them to hold simultaneously. If each individual statement
may fail with probability $\delta$, then over $N$ time steps the chance that \emph{some}
statement fails grows with $N$. So $\delta$ must shrink as $N$ grows. The standard choice
(Auer, Cesa-Bianchi and Fischer, 2002) is
\[
\delta = N^{-4}, \qquad\text{giving}\qquad \ln(1/\delta) = 4\ln N .
\]
Substituting into Equation~\eqref{eq:epsilon} with $R$ and the constants folded into a single
tunable $c$:
\[
\epsilon_i = R\sqrt{\frac{4 \ln N}{2 n_i}} \;=\; \underbrace{R\sqrt{2}}_{\text{call it } c}\,
\sqrt{\frac{\ln N}{n_i}} .
\]

\begin{finalform}[UCB1]
\begin{equation}
  \label{eq:ucb1}
  \boxed{\;\text{pull}\;\; \argmax_{i}\left[\; Q_i \;+\; c\,\sqrt{\frac{\ln N}{n_i}}\;\right]\;}
\end{equation}
with $N$ the total pulls so far, $n_i$ the pulls of machine $i$, and $c > 0$ an exploration
constant. Any machine with $n_i = 0$ is treated as having infinite bonus, so every machine
is pulled once before any is pulled twice.
\end{finalform}

Now read the two moving parts against the failures they repair:
\begin{itemize}
\item $\dfrac{1}{\sqrt{n_i}}$ --- \emph{"I have looked at this $n_i$ times, so my error bar
  is this wide."} Repairs Attempt~1. It is the law of Chapter~\ref{ch:onemachine}, unchanged.
\item $\sqrt{\ln N}$ --- \emph{"...but the search has moved on, and being unexamined is more
  suspicious now than it was."} Repairs Attempts~2 and~3. It grows without bound, so no
  machine is frozen out forever; and it grows only \emph{logarithmically}, so the revival is
  slow and cheap. Doubling the search multiplies this term by $\sqrt{\ln 2N/\ln N}$ --- a
  few percent, not a factor.
\end{itemize}

\begin{keyidea}
The bonus is a \textbf{plausibility of being underrated}. It is large when a machine is
poorly measured relative to how much total measuring has gone on, and it decays as you look.
The rule is not "explore sometimes"; it is "always take the option that could be best, given
what you actually know" --- and that single rule produces exploration as a consequence,
targeted exactly where the ignorance is.
\end{keyidea}

\section{UCB1 by hand}

\begin{worked}{twelve pulls, three machines}
Truth (hidden from the algorithm): $\mu_A = 0.7$, $\mu_B = 0.5$, $\mu_C = 0.2$. To keep the
arithmetic clean, assume each pull of a machine returns exactly its true value plus the
listed noise. Take $c = 1$. \emph{Illustrative numbers, not engine output.}

\smallskip
Pulls 1--3 (each machine once, since $n_i = 0$ gives infinite bonus):
$x_A = 0.3$ (unlucky), $x_B = 0.5$, $x_C = 0.2$. So $Q_A = 0.3$, $Q_B = 0.5$, $Q_C = 0.2$,
all $n_i = 1$, and $N = 3$.

\smallskip
Pull 4: $\ln N = \ln 3 = 1.099$, so every bonus is $\sqrt{1.099/1} = 1.048$.
\[
\text{scores: } A = 1.348,\quad B = 1.548,\quad C = 1.248 \;\Rightarrow\; \text{pull } B.
\]
Say $x_B = 0.5$, so $Q_B = 0.5$, $n_B = 2$, $N = 4$.

\smallskip
Pull 5: $\ln 4 = 1.386$.
\[
A: 0.3 + \sqrt{1.386/1} = 0.3+1.177 = 1.477,\quad
B: 0.5 + \sqrt{1.386/2} = 0.5+0.833 = 1.333,\quad
C: 0.2+1.177 = 1.377 .
\]
$\Rightarrow$ pull $A$ --- \textbf{note what just happened}: $A$ has the \emph{worst} average
of the two leaders, and it gets pulled anyway, because $B$'s bonus shrank when $B$ was
measured. This is the greedy catastrophe being averted, on schedule and without randomness.
Say $x_A = 0.9$, so $Q_A = (0.3+0.9)/2 = 0.6$, $n_A = 2$, $N = 5$.

\smallskip
Pull 6: $\ln 5 = 1.609$; bonuses $\sqrt{1.609/2} = 0.897$ for $A$ and $B$, $\sqrt{1.609/1} = 1.269$ for $C$.
\[
A: 0.6+0.897 = 1.497,\quad B: 0.5+0.897 = 1.397,\quad C: 0.2+1.269 = 1.469 \;\Rightarrow\;
\text{pull } A .
\]
Say $x_A = 0.7$: $Q_A = 0.733$, $n_A = 3$, $N = 6$.

\smallskip
Continuing the same arithmetic to $N = 12$ gives roughly $n_A = 6$, $n_B = 4$, $n_C = 2$
(exact counts depend on the noise you assume; Exercise~\ref{ex:ch4:continue} asks you to
finish it). The shape is the thing:

\begin{center}
\begin{tabular}{@{}lrrr@{}}
\toprule
Machine & true $\mu$ & pulls $n_i$ & share of budget \\
\midrule
$A$ & 0.7 & 6 & 50\% \\
$B$ & 0.5 & 4 & 33\% \\
$C$ & 0.2 & 2 & 17\% \\
\bottomrule
\end{tabular}
\end{center}

The budget is \emph{graded}, not switched. The worst machine is not banned --- it keeps a
trickle --- and the best gets the most without ever being declared the winner. And regret by
Equation~\eqref{eq:regret2} is $4(0.2) + 2(0.5) = 1.8$, against the greedy disaster of the
earlier example.
\end{worked}

\begin{established}
\textbf{The guarantee (Auer et al., 2002).} UCB1's expected regret after $T$ pulls is
$O\!\left(\sum_{i:\Delta_i>0} \frac{\ln T}{\Delta_i}\right)$: it grows only
\emph{logarithmically} in $T$. Since a strategy that pulls any suboptimal machine a constant
fraction of the time has regret growing \emph{linearly}, this is a genuine and strong
result: in the long run UCB1 spends a vanishing fraction of its budget on inferior options,
without ever needing to know $\mu$.
\end{established}

Notice also what the $1/\Delta_i$ says: machines that are \emph{nearly} as good as the best
get pulled a lot. That is not a bug. If two moves are within $0.02$ of each other it does
not much matter which you play, and the search correctly declines to spend its budget
separating them --- while a move that is $0.9$ worse is dismissed in a handful of visits.
Engines look sharp-eyed about blunders and indecisive about near-equal options for exactly
this reason.

\section{The exploration constant}

$c$ is the one free knob, and it does what you would expect:

\begin{center}
\begin{tabular}{@{}ll@{}}
\toprule
$c = 0$ & pure greed; the catastrophe of §4.2 \\
$c$ small & narrow, deep, confident search; risks locking onto an early mistake \\
$c$ large & broad, shallow search; every move gets attention, nothing is resolved \\
$c \to \infty$ & uniform sampling; you have re-invented brute force \\
\bottomrule
\end{tabular}
\end{center}

\begin{realdata}[title={Real engine setting}]
Leela's equivalent constant is exposed as the UCI option \texttt{CPuct}, and on the engine in
\texttt{engine/} its default value is
\[
  \texttt{CPuct} = 1.745,
\]
alongside \texttt{CPuctBase} $= 38739$ and \texttt{CPuctFactor} $= 3.894$, which make it grow
slowly with the size of the search (Chapter~\ref{ch:puct} explains why). Verified by running
\texttt{lc0.exe} with the \texttt{uci} command and reading the option list.
\end{realdata}

\section{Where the guarantee stops}

Be clear about what has and has not been established, because this is exactly the sort of
place where a confident-sounding book does damage.

\begin{pitfall}
The regret theorem assumes independent, identically distributed samples from a fixed
distribution. In a chess search:
\begin{itemize}
\item samples are \emph{not} independent --- successive visits to a move deliberately go
  down different, related lines;
\item the distribution is \emph{not} fixed --- as the subtree below a move grows, the value
  coming back changes systematically (that is the entire point of searching deeper);
\item and the goal is different: we do not want to maximise the sum of payouts along the
  way, we want to \emph{identify the best move at the end}. Those are different objectives
  and are known to have different optimal strategies.
\end{itemize}
So UCB1 is the \textbf{motivation} for what engines do, not a proof that it is optimal for
them. What survives contact with practice is the \emph{shape}: measured value plus a bonus
that shrinks with evidence and grows with total effort. Every strong Monte-Carlo engine uses
that shape; none of them satisfies the theorem's hypotheses.
\end{pitfall}

That honest caveat also opens the door to the next two chapters. Because the theorem does not
bind us, we are free to change the bonus into something better suited to chess --- and in
Chapter~\ref{ch:puct} we will replace the purely statistical "how much have I looked?" with
"how much have I looked, \emph{and how good does a strong player think this move is?}"

\begin{exercises}

\exhead[regret]{ch4:regret} Three moves have true values $0.8, 0.75, 0.1$. A search of 800
visits allocates them $600, 100, 100$. Compute the regret using Equation~\eqref{eq:regret2}.
Now compute the regret for the allocation $400, 390, 10$. Which search would you rather have,
and does regret agree with you?

\exhead[greedy]{ch4:greedy} Construct a two-machine example in which greedy achieves
\emph{zero} regret, and one in which it achieves the maximum possible. What property of the
first sample decides which of these happens? What does this tell you about relying on a
search's first impression of a move?

\exhead[$\epsilon$-greedy's tax]{ch4:epsgreedy} With $K = 31$ moves, $\epsilon = 0.1$, and
$T = 800$ visits, how many visits go to uniformly random moves? If 25 of the 31 moves are
blunders with $\Delta_i \approx 1.0$, roughly how much regret does the random component alone
generate? Compare with the UCB1 behaviour described in §4.4.

\exhead[why $\ln N$ and not $N$]{ch4:lnN} Suppose someone proposes bonus
$= c\sqrt{N/n_i}$ (no logarithm). Take a machine with $n_i = 1$ and $Q_i = -0.9$ and a leader
with $Q = 0.8$, $n = 100$, and compute both scores at $N = 100$ and at $N = 10{,}000$. What
goes wrong, and what does the logarithm buy?

\exhead[bonus table]{ch4:bonustable} Tabulate the UCB1 bonus $\sqrt{\ln N / n_i}$ with $c=1$
for $n_i \in \{1, 4, 16, 64\}$ at $N = 100$ and again at $N = 10{,}000$. By what factor did
each bonus grow when the search got 100 times bigger? State the general rule in one sentence.

\exhead[finish the simulation]{ch4:continue} Continue the worked example from pull 7 to pull
12, assuming every pull of $A$ returns $0.7$, of $B$ returns $0.5$, of $C$ returns $0.2$
exactly. Give the full table of $Q_i$, $n_i$ and scores at each step, and the final visit
distribution.

\exhead[the near-equal case]{ch4:neartie} Two moves have $\mu_A = 0.60$ and $\mu_B = 0.58$.
Using the $\ln T/\Delta_i$ shape of the regret bound, argue qualitatively how the visits
split, and explain why an engine that "cannot make up its mind" between two moves is
usually telling you something true about the position rather than failing.

\exhead[budget and $c$]{ch4:tuningc} You run a search with 200 nodes and one with 20{,}000
nodes. Should $c$ be the same in both? Argue both sides, then look at the real Leela defaults
in §4.6 and say which side the engine's designers took (\texttt{CPuctBase} and
\texttt{CPuctFactor} are the clue; you will confirm your reading in Chapter~\ref{ch:puct}).

\exhead[applying it to chess, carefully]{ch4:tochess} List three specific ways in which
visiting a chess move differs from pulling a slot machine. For each, say whether it makes
UCB1's assumption more or less reasonable, and name (guessing is fine) what an engine might
do to compensate.

\exhead[the diagnostic]{ch4:diagnostic} Your tool reports that in one of your games, the
engine spent $92\%$ of its visits on a single move. In another position it spread visits
$30/28/25/17$ across four moves. Using this chapter's vocabulary only --- no chess ---
describe what each distribution says about the position. (Chapter~\ref{ch:readingtree} turns
your answer into a feature.)

\end{exercises}
